\chapter{Aleaciones metálicas aeroespaciales}\label{ch:07}
\objetivos{Comparar familias de aleaciones sin convertir una denominación comercial en una especificación; relacionar fases, tratamientos y aplicaciones.}
\section{Cómo leer una designación}
La composición nominal identifica una familia, pero la propiedad final exige conocer producto, espesor, proceso, tratamiento y dirección. En aluminio, la serie describe el sistema principal de aleación; el estado de suministro añade información esencial. En titanio se consideran las fases y el tratamiento. En superaleaciones importa también si el producto es forjado, colado, de solidificación direccional, monocristalino o de pulvimetalurgia.

La selección necesita especificaciones y datos admisibles vinculados al producto. Una tabla académica puede organizar tendencias; no constituye una autorización para sustituir una aleación por otra «equivalente» en un motor.
\section{Aceros aeronáuticos e inoxidables}
Los aceros ofrecen combinaciones amplias de resistencia, tenacidad, endurecimiento y coste. En aceros templados y revenidos se ajustan propiedades mediante transformaciones y tratamiento posterior. Los aceros maraging obtienen gran parte del endurecimiento de precipitación en una matriz martensítica de bajo carbono; no son simplemente un acero al carbono con mayor temple.

Los inoxidables se apoyan en la pasividad asociada al cromo, con familias austeníticas, ferríticas, martensíticas y endurecibles por precipitación, entre otras. Su resistencia a corrosión no significa inmunidad frente a picadura o corrosión bajo tensión. La elección de familia depende también de temperatura, magnetismo, soldabilidad y resistencia mecánica.

En ejes, engranajes, elementos de unión y rodamientos importan resistencia a contacto, desgaste, tenacidad y estabilidad dimensional. La alta resistencia puede aumentar sensibilidad a defectos o fragilización, de modo que el tratamiento superficial y el control del proceso no son detalles secundarios.
\section{Aleaciones de aluminio}
Las series 2xxx se asocian principalmente con cobre, las 6xxx con magnesio y silicio, y las 7xxx con zinc como elemento principal, con otros elementos que modifican las fases de endurecimiento. Estas familias incluyen aleaciones tratables térmicamente, pero no todas comparten propiedades ni resistencia a corrosión.

Los estados T3, T6 o T7 describen historias diferentes. El aumento de resistencia por maduración puede implicar compromisos en tenacidad o susceptibilidad ambiental. Un tratamiento sobreenvejecido puede buscar resistencia a corrosión bajo tensión a costa de parte de la resistencia máxima; el compromiso debe justificarse con datos del producto.

Las aleaciones aluminio--litio permiten reducir densidad y modificar rigidez específica, pero plantean cuestiones de anisotropía, tenacidad, fabricación y uniones. La designación 8090 corresponde a una aleación de aluminio--litio; no es una aleación de níquel. Tampoco una temperatura de servicio inferior a $200\,^{\circ}\mathrm C$ descarta «todos los metales»: la selección depende de los requisitos y de la estabilidad de cada aleación.
\section{Titanio: fases alfa y beta}
El titanio presenta fase alfa de estructura hexagonal compacta y beta de estructura cúbica centrada en el cuerpo. Los aleantes estabilizadores y el procesamiento controlan su proporción y morfología. Se distinguen aleaciones alfa, cercanas a alfa, alfa--beta y beta, con respuestas diferentes al tratamiento.

El titanio combina baja densidad relativa frente a aceros y superaleaciones con resistencia específica y buen comportamiento en ambientes adecuados. Sus límites incluyen coste, dificultad de mecanizado, sensibilidad del proceso y compatibilidad a alta temperatura. La interacción con oxígeno durante operaciones calientes puede alterar la superficie; no basta con indicar una resistencia elevada a temperatura ambiente.

En compresores y elementos del fan puede ser atractivo, pero la selección debe considerar temperatura, vibración, impacto, desgaste por contacto y riesgos de ignición o compatibilidad con materiales vecinos. No se deduce una temperatura máxima universal de la palabra «titanio».
\section{Superaleaciones de níquel y cobalto}
En muchas superaleaciones de níquel, una matriz gamma contiene precipitados gamma prima, habitualmente basados en $\mathrm{Ni_3(Al,Ti)}$, que obstaculizan deformación. Otras familias emplean mecanismos y fases diferentes, como gamma doble prima ($\gamma^{\prime\prime}$) en determinadas aleaciones. Carburo, composición de límites y tratamiento influyen en la respuesta. Una fracción de precipitados favorable a una temperatura puede evolucionar durante exposición prolongada.

La resistencia a fluencia, oxidación y fatiga termomecánica hace útiles estas familias en zonas calientes. La densidad y el coste son altos; el diseño de discos y álabes no tiene idénticas prioridades. Las superaleaciones de cobalto pueden ofrecer prestaciones térmicas, ambientales o de desgaste apropiadas en aplicaciones específicas, pero no son un reemplazo universal de las de níquel.

Los álabes monocristalinos eliminan límites de grano de alto ángulo dentro del cristal, aunque mantienen anisotropía y defectos posibles de fabricación. Los discos de grano fino o estructura controlada priorizan otros equilibrios entre fatiga, resistencia y tenacidad. La relación entre microestructura y misión se desarrolla en el capítulo~\ref{ch:10}.
\begin{table}[htbp]\centering\small\caption{Comparación cualitativa de familias. Las prestaciones concretas requieren datos del producto.}\label{tab:familias}
\begin{tabular}{p{.18\linewidth}p{.33\linewidth}p{.37\linewidth}}\toprule
Familia & Ventajas que pueden motivar la selección & Aspectos que deben comprobarse\\\midrule
Aceros & Resistencia, desgaste y diversidad de tratamientos & Masa, corrosión, tenacidad y fragilización\\[4pt]
Aluminio & Densidad baja, fabricación y resistencia específica & Temperatura, fatiga, uniones y ambiente\\[4pt]
Titanio & Resistencia específica y compatibilidad ambiental & Coste, proceso, contacto y límite térmico\\[4pt]
Níquel/cobalto & Servicio caliente y estabilidad microestructural & Densidad, fluencia, recubrimientos y procesado\\\bottomrule
\end{tabular}\end{table}
\section{Una selección explicada}
\caso{M07. Disco y álabe no piden lo mismo}{Se comparan dos piezas ficticias a distinta temperatura: un disco sometido a ciclos centrífugos y un álabe sometido a un gradiente térmico con permanencias. El análisis comienza por fatiga y tolerancia a defectos en el disco, y por fluencia, oxidación y fatiga térmica en el álabe, sin excluir los demás modos. Por ello, escoger la misma familia química no obliga a usar el mismo tamaño de grano ni el mismo proceso. La conclusión es una estrategia de comparación, no una especificación de fabricación.}
\section{Identificación de materiales sin respuestas absolutas}
La aleación 8090 es aluminio--litio, no níquel. Las fases del titanio son alfa, de red hexagonal compacta (HCP), y beta, cúbica centrada en el cuerpo (BCC). Ti--6Al--4V puede presentar microestructuras alfa+beta según proceso y tratamiento: el nombre comercial no permite fijar una microestructura única para cualquier pieza. Los aceros maraging se endurecen principalmente por precipitación, lo cual no garantiza por sí solo resistencia universal a la corrosión o tenacidad suficiente.

Una pregunta sobre material de tren de aterrizaje, fan o estructura de aeronave no admite elegir sin conocer modelo, zona, requisitos y estado metalúrgico. Deben considerarse fatiga, fractura, corrosión, fabricación, temperatura e inspección. La temperatura de fusión de una superaleación no determina por sí sola su máxima temperatura admisible de servicio.


\section{Comparar familias sin convertirlas en recetas}
En aluminio, las series 2xxx y 7xxx suelen relacionarse con endurecimiento por precipitación y elevadas resistencias en estados concretos; las 6xxx ofrecen otras combinaciones de fabricación y propiedades. La familia no determina una resistencia única ni garantiza soldabilidad, resistencia a corrosión o tenacidad. Un estado de tratamiento que maximiza una propiedad puede penalizar otra. Las aleaciones Al--Li añaden posibilidades de menor densidad y modificación del módulo, junto a exigencias de procesado y comportamiento que deben evaluarse en el producto.

En titanio, las proporciones y morfologías de fases alfa y beta dependen de composición e historia térmica. Las propiedades no se deducen únicamente del nombre de la fase: tamaño, distribución y textura importan. Se valoran resistencia específica y compatibilidad térmica, pero también fabricación, comportamiento de superficies y posibles daños por contacto. No se interpreta «resistente a corrosión» como inmunidad en todo ambiente.
\section{Discos y álabes: diferentes compromisos}
Un disco debe soportar esfuerzos asociados a rotación, gradientes térmicos, transitorios y concentraciones, con una población de defectos y una vida controladas. Un álabe caliente añade exposición severa a temperatura, ambiente y vibración y puede usar refrigeración y recubrimientos. Una microestructura adecuada para uno no es automáticamente la mejor para el otro.

Las superaleaciones de níquel combinan matriz y fases de refuerzo estables en un intervalo de servicio; composición, fracción, morfología y estabilidad de precipitados influyen en fluencia y fatiga. El procesado de disco o álabe busca microestructuras diferentes. Las superaleaciones de cobalto también ofrecen combinaciones útiles en aplicaciones calientes, pero no se sustituyen por las de níquel usando solo la temperatura de fusión. El comportamiento de oxidación, corrosión, fabricación y unión completa la comparación.
\section{Aceros, inoxidables y coste de ciclo de vida}
Aceros de alta resistencia pueden resultar apropiados para ejes, engranajes o elementos muy cargados, según diseño. Esa resistencia puede acompañarse de mayor sensibilidad a entallas, ambiente o tratamientos. Los inoxidables reúnen familias microestructurales distintas; la pasividad depende del medio y de la condición superficial. No existe un inoxidable universalmente inmune a picaduras o corrosión bajo tensión.

El coste de adquisición es solo una parte: herramientas, inspección, reparación, disponibilidad de producto, tiempos de suministro y retirada de servicio afectan al coste de ciclo de vida. Una preselección técnicamente prometedora puede descartarse si no existe una ruta de fabricación y control reproducible para la geometría requerida.
\repaso{¿La designación 8090 identifica una superaleación de níquel?}{No: corresponde a una aleación de aluminio--litio. La identificación de familia se verifica antes de discutir propiedades; un número de designación aislado no basta para asignar estado o admisibles.}

\resumen{La designación de aleación es solo parte del dato. Las series de aluminio no sustituyen el estado metalúrgico. Las fases del titanio y los precipitados de las superaleaciones dependen del procesamiento. Temperatura, daño, inspección, fabricación y coste deben compararse conjuntamente.}
\fuentes{FAA--General, cap. 7 \cite{FAA-G}; NASA, desarrollo de superaleaciones de baja densidad \cite{NASA-SX}; materiales docentes de introducción. Para valores de diseño se necesitaría la especificación y su base estadística aplicable.}
